\subsection{Centering the influence function before weighting and clustering}
\label{rel23:ife-centering}

The selected uncertainty delegate sharpens the distinction between a policy
value and the uncertainty of that value. The inspected \texttt{ife} revision
declares version 0.2.5 in July 2026 and GPL version 3 or later. It stores a
point estimate, a vector called the influence function, weights, cluster
identifiers and a critical value. Arithmetic operations propagate derivatives
through the vector. This is a useful interface for retaining the dependence
between two policies evaluated on the same observations. It does not construct
the correct influence function from an arbitrary prediction vector. The
separately inspected lmtp revision permits \texttt{ife} versions from 0.2.0;
its resolved dependency is not observed here. The delegate's own change log
records changes to cluster handling in 0.2.1 and 0.2.5, so the lower version
bound cannot identify the behavior of a particular installation.
\citep{rel23ife,exec22_lmtp_reuse}

The boundary matters because the retained lmtp caller supplies an
\emph{uncentered} sequential pseudo-outcome: weighted residual contributions
plus an initial shifted prediction. It rescales that vector and its point
estimate to the outcome's original units, then forwards both to \texttt{ife}.
The delegate stores the supplied vector without subtracting the estimate.
For independent, equally weighted rows, sample variance removes a common
constant. That familiar simplification fails when multiplication by unequal
weights or aggregation over unequal cluster sizes occurs first.
\citep{exec22_lmtp_reuse,rel23ifecaller,rel23ife}

To see the statistical issue without fitting a nuisance model, let $Z_i$ be
a supplied pseudo-outcome and $w_i$ a nonnegative sampling weight. Suppose
the target is the ratio of weighted population moments, estimated by
\begin{equation}
 \widehat\theta=\frac{\sum_i w_iZ_i}{\sum_i w_i},\qquad
 \widetilde w_i=\frac{w_i}{n^{-1}\sum_j w_j},\qquad
 D_i=\widetilde w_i(Z_i-\widehat\theta).
 \label{rel23:weighted-influence}
\end{equation}
Here $n$ is the number of independent observations and $D_i$ is the empirical
influence contribution for this ratio. Differentiating the numerator divided
by the denominator gives the subtraction in the last expression; it accounts
for variation in the denominator as well as the weighted outcome. With
positive finite mean weight, finite second moments and an appropriate
independent sampling law, the usual first-order variance estimate is the
sample variance of these contributions divided by $n$. Fitted longitudinal
pseudo-outcomes additionally require the identifying and nuisance-estimation
conditions developed earlier. This elementary ratio calculation does not
supply those conditions. \citep{rel23ifewitness}

The lmtp task mean-normalizes supplied weights. The selected delegate,
however, first multiplies its supplied vector by those weights. In an exact
authored example, take pseudo-outcomes $(1,2,4,5)$ dollars and normalized
weights $(1/2,1/2,3/2,3/2)$. The weighted mean is USD~3.75. Applying the
displayed centered formula gives a variance estimate of $101/192$ squared
dollars and standard error USD~0.725. The uncentered weighted vector instead
gives $149/48$ and standard error USD~1.762 under the documented sample-variance
convention. Add a known USD~100 accounting offset to every pseudo-outcome.
The centered contributions are unchanged, as uncertainty about an unknown
mean should be under a known shift. The uncentered formula now gives
$14983/16$ and standard error USD~30.601. These are exact source-formula
calculations, not executed R results or estimates of actual bank uncertainty.
\citep{rel23ife,rel23ifecaller,rel23ifewitness}

The same defect can survive a contrast. Consider a stipulated second reward
equal to the first reward plus a known USD~2 for every account, using the
same weights. Their estimated weighted-mean difference is identically two;
its influence contribution is zero. Subtracting the two uncentered vectors
instead yields a constant two, which becomes the unequal vector of twice
the weights when the delegate calculates uncertainty. Its sample-variance
estimate divided by four is $1/3$ squared dollars. This example is a known
reward transformation, not evidence of a constant causal policy effect. It
shows why subtracting point estimates and retaining paired rows still does
not repair an incorrectly centered uncertainty representation.
\citep{rel23ifewitness}

For several independent clusters, the inspected delegate sums weighted
contributions within each cluster and multiplies each sum by the number of
clusters divided by the number of rows. With correctly centered individual
contributions, this gives the empirical cluster influence
\begin{equation}
 G_k=\frac{K}{n}\sum_{i\in C_k}D_i,\qquad
 \widehat{\operatorname{Var}}(\widehat\theta)
       =\frac{s_G^2}{K},
 \label{rel23:cluster-influence}
\end{equation}
where $C_k$ is cluster $k$, $K$ is the number of independent clusters and
$s_G^2$ is their sample variance. Cluster-level inference needs a justified
sampling law and enough independent clusters without a dominating cluster.
The target here remains an individual-weighted mean; averaging branch means
with equal branch weight would define a different target. The presence of an ID column
cannot provide them. As an algebraic check, use four unweighted outcomes
$(0,0,2,2)$, with the first account in one cluster and the remaining three in
another. Centering at their mean gives cluster contributions $(-1/2,1/2)$
and variance estimate $1/4$. The uncentered calculation uses $(0,2)$ and gives
one. Adding the same known USD~100 offset changes these latter values to
$(50,152)$ and the variance estimate to 2,601. Two clusters are sufficient to
expose this origin dependence; they are not a calibration study. The source
also deliberately treats exactly one cluster as independent rows, a choice
its test asserts. An adapter requiring cluster-robust uncertainty should
return that uncertainty as unavailable when only one independent cluster
exists, instead of silently substituting the independent-row sampling law.
\citep{rel23ife,rel23ifewitness}

\subsection{A paired contrast must retain its scale and sampling identities}
\label{rel23:ife-contrasts}

The delegate's arithmetic is useful once the incoming influence representation
is correct. A difference subtracts corresponding contributions; a ratio uses
the numerator derivative and the denominator derivative jointly. This retains
their covariance. Before combining two objects, the inspected code requires
equal vector lengths and identical weight and ID arrays. Those checks help,
but cluster labels do not uniquely identify the individuals within a cluster.
The delivery contract must bind each influence contribution to the same
original observation across fits, as well as to its cluster and weight.
Reordering records while regenerating positional IDs can pass an array-level
check while changing the pairing. Such a change alters uncertainty without
changing either marginal point estimate. \citep{rel23ife}

The newly inspected lmtp contrast path computes relative-risk and odds-ratio
contrasts on the log scale, obtains a symmetric interval there and then
exponentiates its endpoints. It also exponentiates the point estimate, but
leaves the returned \texttt{std.error} column on the log scale. That is usable
when documented; interpreting every numeric column as being on the reported
ratio scale is a conversion error. For a ratio of 1.5 and log-scale standard
error 0.2, the delta-method ratio-scale standard error is 0.3. A conventional
normal critical value gives an exponentiated log interval approximately
$[1.014,2.220]$, compared with $[0.912,2.088]$ for a symmetric interval formed
on the ratio scale. Neither arithmetic construction proves finite-sample
coverage, especially near a weak denominator. The delegate's rejection of
denominators smaller than machine epsilon is a numerical guard, not a test
that a denominator is statistically separated from zero.
\citep{rel23ifecaller,rel23ife,rel23ifewitness}

The caller warns when a numeric reference is used because an estimated
reference treated as known loses its uncertainty. Its print method warns
about unadjusted multiple comparisons. The delegate permits a configurable
critical value but prints a fixed 95\% interval label. A bank-facing result
therefore needs the actual contrast, scale of every uncertainty field,
critical value, sampling unit and comparison family. A method name or a
formatted confidence label should not supply these facts by implication.
\citep{rel23ifecaller,rel23ife}

The recommendation is conditional reuse of the arithmetic interface after
qualifying the incoming centered influence function, mean-weight normalization,
individual pairing and independent clusters. Compare a known outcome shift,
a deterministic reward increment, unequal cluster sizes, a single-cluster
request and log-versus-ratio output against independent laws before relying
on a standard error. A simpler independent implementation of the declared
influence calculation is a serious alternative for a narrow initial worker;
it must meet the same statistical requirements. Native dependency resolution,
the numerical variance delegate, full survival handling, learner behavior and
random-number coordination remain separate qualification obligations. These
source findings strengthen the admission conditions for the longitudinal
reference; they do not establish an executed package failure rate, calibrated
bank uncertainty or an observed business benefit.
\citep{rel23ife,rel23ifecaller,rel23ifewitness}
